% !TeX root = ../深度学习与大语言模型技术原理.tex
% 由主文档通过 \input 载入；请编译主文档。
\setcounter{chapter}{0} % 保留原书章号。
\chapter{技术清单与正文对应索引}
\label{app:map}

表~\ref{tab:source-map} 的左栏使用来源文件本身的节号，便于核对范围；右栏使用本书自动交叉引用。清单中的重复主题合并到同一论述，技术栈名称则放在相应方法或工程章节。清单的优先级与求职描述只用于识别学习主题，不作为本书的经验结论。

\begin{longtable}{L{0.12\textwidth}L{0.45\textwidth}L{0.31\textwidth}}
\caption{来源条目与本书内容的对应关系}\label{tab:source-map}\\
\toprule
清单节号 & 技术主题 & 本书位置 \\
\midrule
\endfirsthead
\toprule
清单节号 & 技术主题 & 本书位置 \\
\midrule
\endhead
\bottomrule
\endfoot
0 & 能力总览、世界模型、稀疏与混合架构 & 第~\ref{ch:introduction}、\ref{ch:moe}、\ref{ch:agentlearning} 章 \\
1 & 数学、信息论、优化 & 第~\ref{ch:probability}、\ref{ch:optimization} 章 \\
2 & 注意力、RoPE、归一化、SwiGLU、GQA、MQA & 第~\ref{ch:transformer} 章 \\
3 & MoE、路由、容量、负载均衡、专家并行 & 第~\ref{ch:moe} 章 \\
4 & 分词、语言建模、预训练数据与配比 & 第~\ref{ch:pretraining}、\ref{ch:data} 章 \\
5 & 混合精度、累积、裁剪、检查点、显存 & 第~\ref{ch:stability} 章 \\
6 & DDP、FSDP、ZeRO、TP、PP、CP、EP & 第~\ref{ch:distributed} 章 \\
7--8 & LoRA、QLoRA、SFT、掩码与拼接 & 第~\ref{ch:sft} 章 \\
9 & 指令数据清洗、合成与筛选 & 第~\ref{ch:data}、\ref{ch:distillation} 章 \\
10 & 奖励模型、偏好与 RLHF & 第~\ref{ch:preference}、\ref{ch:policy} 章 \\
11--12 & PPO 与 DPO & 第~\ref{ch:policy}、\ref{ch:preference} 章 \\
13 & IPO、KTO、ORPO、SimPO、RLOO & 第~\ref{ch:preference}、\ref{ch:policy} 章 \\
14--16 & 推理后训练、GRPO 与 RLVR & 第~\ref{ch:policy} 章 \\
17--18 & 结果奖励、过程奖励与验证器 & 第~\ref{ch:policy}、\ref{ch:evaluation} 章 \\
19--21 & 拒绝筛选、推理蒸馏、推理时计算 & 第~\ref{ch:distillation} 章 \\
22--24 & 智能体循环、工具调用、MCP & 第~\ref{ch:tools}、\ref{ch:agents} 章 \\
25--28 & 记忆、上下文、规划、多智能体 & 第~\ref{ch:agents} 章 \\
29--30 & 智能体强化学习、计算机操作 & 第~\ref{ch:agentlearning} 章 \\
31 & RAG 的整体模型与流程 & 第~\ref{ch:rag} 章 \\
32--35 & 稀疏、稠密、混合检索与重排 & 第~\ref{ch:retrieval} 章 \\
36--39 & 查询变换、上下文化、迭代与图检索 & 第~\ref{ch:rag} 章 \\
40 & 检索和生成的分层评测 & 第~\ref{ch:rag} 章 \\
41--45 & 基准、裁判、评分器、在线评测 & 第~\ref{ch:evaluation} 章 \\
46--47 & 自回归推理、预填充、解码与 KV & 第~\ref{ch:inference} 章 \\
48--52 & 分页、连续批处理、前缀缓存、分块、FlashAttention & 第~\ref{ch:efficientattention} 章 \\
53--54 & 权重、激活与 KV 量化 & 第~\ref{ch:quantization} 章 \\
55--57 & 推测解码、分离服务与服务框架 & 第~\ref{ch:speculation} 章 \\
58 & 蒸馏、量化、剪枝 & 第~\ref{ch:distillation}、\ref{ch:quantization} 章 \\
59 & 长上下文、位置扩展、训练与评测 & 第~\ref{ch:longcontext} 章 \\
60 & 结构化输出与约束解码 & 第~\ref{ch:tools} 章 \\
61 & 图文、音频、视频与统一多模态 & 第~\ref{ch:multimodal} 章 \\
62 & 代码预训练、填充、仓库与软件智能体 & 第~\ref{ch:agentlearning} 章 \\
63--65 & 提示注入、工具护栏与沙箱 & 第~\ref{ch:safety} 章 \\
66 & 探针、激活替换、因果追踪、SAE & 第~\ref{ch:interpretability} 章 \\
67--69 & GPU、CUDA、Triton、编译与 NCCL & 第~\ref{ch:kernels}、\ref{ch:distributed} 章 \\
70--71 & Python、C++、Linux、Git、Docker 与工具链 & 第~\ref{ch:engineering} 章及相应系统章节 \\
72 & 重复出现的前沿主题 & 按各技术所在章节查阅 \\
73--74 & 分方向学习与先修次序 & 第~\ref{ch:introduction} 章及本附录后文 \\
75 & 六个综合项目 & 附录~\ref{app:experiments} \\
76--77 & 原理解释、实现、优化与能力检验 & 正文推导及附录~\ref{app:experiments} \\
78--80 & 模型、数据、评测和系统的完整流程 & 第~\ref{ch:engineering} 章 \\
81--82 & 原始资料与持续学习 & 书末参考文献及本附录后文 \\
\end{longtable}

\section{按研究问题选择深化方向}

后训练方向应把 SFT、奖励、策略梯度、采样系统与评测连起来学。能写出 DPO 损失只是起点，还要判断偏好数据是否可靠、GRPO 奖励是否有信息、生成吞吐是否成为瓶颈，以及训练改善是否转化为独立任务收益。检索方向则需要同时理解表示学习、索引误差、排序和证据生成，不能把所有失败都归于向量数据库。

智能体方向应先建立状态、工具契约与评测，再研究规划和学习；系统方向应先会核算计算、显存和通信，再选择量化、并行或内核优化。基础模型方向要把数据分布、模型容量和训练预算联合分析。以上是知识依赖建议，不是招聘需求或薪资预测。

持续阅读可以从正文中遇到的具体假设出发，回到原始论文的目标、实验设置和消融。新方法出现时，先判断它改变了目标、数据、架构还是执行方式，再检查是否解决当前问题。论文复现应固定版本并保留负结果；参考文献中的动态文档是入口，实际运行仍需记录所用版本。

